\documentclass[12pt]{article}
\usepackage[margin=0.75in]{geometry}

\newcommand{\duedate}{04/05/2026}
\newcommand{\assignment}{End-to-End Test-Time Training for Long Context: Summary}

\newcommand{\name}{Aksel Kretsinger-Walters, adk2164}
\newcommand{\email}{adk2164@columbia.edu}

\makeatletter
\def\input@path{{../}{../../}{../../../}}
\makeatother
\input{pset_template_CLMM.tex} %% DO NOT CHANGE THIS LINE

\begin{document}
\psetheader %% DO NOT CHANGE THIS LINE

\section*{End-to-End Test-Time Training for Long Context}

\paragraph{Motivation.}
Transformers with full self-attention achieve strong long-context performance but
have $O(T^2)$ prefill cost and $O(T)$ per-token decode cost, which becomes
prohibitive for very long sequences. RNN-style alternatives (Mamba, DeltaNet)
offer constant cost per token but empirically degrade as context grows.
Sliding-window attention (SWA) provides a middle ground but cannot access
information beyond the window.

The key insight of this paper is that long-context language modeling can be
framed as a \emph{continual learning} problem rather than an architecture
design problem.

\paragraph{Core Idea: TTT via Next-Token Prediction.}
Instead of designing new architectures to store long-range context, the authors
propose continuing to \emph{train} a standard Transformer (with SWA) at test
time on the context itself, using the same next-token prediction loss used
during pretraining. The model compresses context into its MLP weights through
gradient descent:
\[
		W_t = W_{t-1} - \eta \, \nabla \ell_t(W_{t-1}),
		\quad \text{where} \quad
		\ell_t(W) = \mathrm{CE}\bigl(f(x_{t-1}; W),\, x_t\bigr).
\]

After processing $T$ tokens of context, the updated weights $W_T$ encode
information beyond the sliding window, enabling the model to attend to the
full context implicitly.

\paragraph{End-to-End (E2E) Training.}
Prior TTT methods (e.g., TTT-KVB) use layer-wise reconstruction losses at
training time but next-token prediction at test time, creating a mismatch.
TTT-E2E resolves this by making the method end-to-end in two senses:

\begin{itemize}
		\item \textbf{E2E at test time:} The inner loop directly optimizes the
				next-token prediction loss (not a layer-wise proxy).
		\item \textbf{E2E at training time:} The outer loop optimizes $W_0$ via
				meta-learning so that the \emph{post-TTT} loss is minimized, using
				gradients of gradients through the inner loop.
\end{itemize}

The training-time objective is:
\[
		\mathcal{L}(W_0; X) = \frac{1}{T} \sum_{t=1}^{T} \ell_t(W_{t-1}),
\]
where each $W_t$ depends on $W_0$ through the chain of gradient updates.

\paragraph{Mini-Batch TTT and Sliding Window.}
Processing tokens one at a time is sequential and unstable. TTT-E2E groups
tokens into mini-batches of size $b$, taking one gradient step per batch.
This enables parallelism within each batch and improves stability. The model
uses SWA layers for local context (window size $k$) and TTT for beyond-window
context, with $k \geq b$ so the window covers at least one full TTT batch.

\paragraph{Implementation Details.}
Three design choices are critical:

\begin{enumerate}
		\item \textbf{TTT only the MLP layers:} Updating attention or normalization
				layers in the inner loop causes instability in the outer loop.
		\item \textbf{TTT only the last $L/4$ blocks:} Updating fewer layers reduces
				compute while maintaining a large effective hidden state. The last
				layers are chosen as they benefit most from context compression.
		\item \textbf{Two MLPs per block:} In TTT-updated blocks, a second frozen
				MLP preserves pretrained knowledge, preventing catastrophic forgetting
				of the base model's capabilities during test-time adaptation.
\end{enumerate}

\paragraph{Connection to Prior TTT Work.}
The paper provides an alternative derivation starting from TTT-KVB (Key-Value
Binding). TTT-KVB uses per-layer reconstruction losses and separate key/value
projections ($\theta_K, \theta_V$). TTT-E2E simplifies this by: (1) removing
the separate output rule, (2) replacing layer-wise losses with the global
next-token prediction loss, and (3) updating fewer but larger MLP layers.
This yields a $5\times$ larger effective hidden state with half the latency.

\paragraph{Scaling Results.}
For 3B-parameter models trained on 164B tokens:

\begin{itemize}
		\item TTT-E2E scales with context length comparably to full attention ---
				both maintain or improve loss from 8K to 128K tokens.
		\item Mamba~2, Gated DeltaNet, and TTT-KVB all degrade relative to full
				attention as context grows beyond 32K.
		\item TTT-E2E has constant inference latency regardless of context length
				(like RNNs/SWA), making it $2.7\times$ faster than full attention
				at 128K context.
\end{itemize}

Results hold across model sizes from 125M to 3B, with consistent improvements
from TTT-E2E on top of SWA even when SWA already has access to full attention
at 8K context.

\paragraph{Limitations.}
TTT-E2E requires backpropagation through the inner-loop gradient steps during
training, increasing training cost. The method currently updates only MLP layers
in the last quarter of the network, leaving open whether more aggressive
updating could help at even longer contexts. The meta-learning outer loop
adds complexity compared to standard pretraining.

\paragraph{Takeaway.}
TTT-E2E reframes long-context modeling as continual learning: rather than
building architectures with perfect recall, compress the context into model
weights via test-time gradient descent. By making both the inner and outer
loops end-to-end with next-token prediction, and leveraging meta-learning to
prepare the initialization, TTT-E2E matches full-attention quality at
RNN-like inference cost.

\end{document}
